\section{Simon of Bethsaida}

\subsection{Three names, and what each carries}

The sources give him three names, and each carries information. His given
name was Simon; the Fourth Gospel adds his patronymic---``Thou art Simon the
son of Jona'' (Jn 1:42), while its final chapter addresses him as ``Simon son
of John'' (Jn 21:15--17). That variation is not a slip in the Douay--Rheims.
It is a genuine divergence in the Greek transmission, and it can be seen in
the apparatus of any edition that records one: the Byzantine textform reads
\textit{Iona} at both Jn 1:42 and Jn 21:15--17, and notes at each of those
verses that the Nestle--Aland text reads \textit{Ioannou} instead. Matthew's
blessing preserves a third form again, the Aramaic patronymic
\textit{Bar-Jona} (Mt 16:17). The father's name is therefore transmitted in
two forms that the surviving witnesses do not reconcile, and no argument in
this study depends on which is original.

His second name was conferred: ``thou shalt be called Cephas, which is
interpreted Peter'' (Jn 1:42). \textit{Cephas} transliterates the Aramaic
word for rock; it is carried untranslated into Paul's Greek letters, where it
is the apostle's usual name (\textit{Kephas} at 1~Cor 1:12, 9:5, 15:5,
Gal 1:18, 2:9), and \textit{Petros} is its Greek rendering, which the Gospels
and Acts prefer and which every later language inherited. Two things follow
that are worth stating plainly at the outset. First, the double naming is
attested independently in Matthew, Mark, and John, though each places it
differently, and only Matthew explains it. Second, John's clause ``which is
interpreted Peter'' is itself evidence that the evangelist expected readers
who did not know Aramaic---the name had already crossed a language boundary
by the time the Gospel was written, and the Greek form had won.

A fourth form is rarer and more revealing. At the Jerusalem council James
introduces his ruling by saying, in the Greek, ``\textit{Symeon} hath related
how God first visited to take of the Gentiles a people to his name''
(Acts 15:14), using the archaic Semitic spelling rather than the ordinary
Greek \textit{Simon}. The only other place in the New Testament where that
spelling is applied to Peter is the opening address of the Second Epistle,
``\textit{Symeon} Peter, servant and apostle of Jesus Christ'' (2~Pt 1:1).
The coincidence is small, and it will be weighed rather than pressed when the
letters are graded below; but it is a real datum, and it is invisible in
translation.

\subsection{The lake and the household}

His world was the fishing economy of the Sea of Galilee. John makes
``Bethsaida, the city of Andrew and Peter'' (Jn 1:44) his place of origin;
the Synoptics locate his working life at Capernaum, where Jesus enters ``the
house of Simon and Andrew'' and heals ``Simon's wife's mother'' lying sick
of a fever (Mk 1:29--31; parallels Mt 8:14--15, Lk 4:38--39). The two
statements are compatible without being one statement: a man may be born in
one lakeside town and keep a household in another, and Bethsaida-origin with
Capernaum-household is this study's synthesis of two witnesses rather than a
claim either of them makes alone.

What the texts show of the trade is modest but consistent. Simon and Andrew
cast a net from the shore; James and John are in a boat mending theirs, and
Mark adds that Zebedee had hired men, so that at least one household in the
partnership employed labour beyond the family (Mk 1:16--20). Luke calls the
sons of Zebedee Simon's \textit{koinonoi}---partners in the ordinary
commercial sense of men who share a venture---and puts Simon in charge of a
boat with a crew large enough to need a second boat's help with a heavy catch
(Lk 5:7, 10). None of this makes him rich; all of it makes him something
other than destitute. He is a working owner-operator with a boat, a net, a
partnership, and a house.

He was married. The Synoptics establish it incidentally, by healing his
mother-in-law, and Paul confirms it decades later in the course of defending
his own rights: have we not the power ``to carry about a woman, a sister, as
well as the rest of the apostles, and the brethren of the Lord, and Cephas?''
(1~Cor 9:5). The Greek phrase is \textit{adelphen gynaika}, literally a
sister-woman or sister-wife; the construction is why translations differ over
whether Paul means a wife or a believing woman who travels with the apostle
as a companion and supporter. What the verse establishes either way is that
in the fifties, when Paul wrote, the churches knew that Cephas travelled
accompanied by a woman, and that this was uncontroversial enough to be used
as a precedent in an argument. Nothing in any first-century source names her,
dates the marriage, or mentions children.

\subsection{Speech, literacy, and the one hostile word}

His speech marked him as Galilean even in Jerusalem: ``even thy speech doth
discover thee'' (Mt 26:73). The only comment any source makes on his
education is the Jerusalem council's astonishment in Acts, that Peter and
John were \textit{agrammatoi} and \textit{idiotai}---``illiterate and
ignorant men'' in the Douay's rendering (Acts 4:13). Both Greek words are
narrower than the English suggests. The first means unlettered in the sense
of lacking formal scribal schooling; the second means a private person, a
non-specialist, as opposed to one trained in a discipline. The sentence is
also placed on the lips of hostile officials inside a narrative written to
show that untrained men spoke with \textit{parrhesia}, boldness. It is
therefore evidence of two things---that Peter had no rabbinic or professional
training, and that Luke wished his readers to notice the contrast---and it is
not a statement about whether he could read, write, or dictate. That
distinction matters later, when the Greek of 1~Peter is weighed, and it is
one of the places where a great deal of confident argument has been built on
two adjectives inside a hostile speech-frame.

\witnesslimit{Nothing else is recoverable. No source records Peter's birth
year, age, appearance, wealth, children, or his wife's name; later
traditions that supply such details (including the legend that made the
Roman martyr Petronilla his daughter) are treated in the tradition audit.
The reconstruction above uses only what the texts say and what the Greek
words mean; it makes no use of archaeology at Capernaum or Bethsaida, no
site was visited, and no excavation report on either town was consulted for
this study.}
